\section{Conclusions}
\label{sec:conclusion}

In this work, a standard for the reproducible assessment of CALPHAD databases, ORCHID-SAGE, has been introduced and tested on the Cr--Ni system. The main idea of the standard is that the record of an assessment has to include not only its data and code, but also the decisions of the assessor, and for this reason an assessment is recorded at three levels, namely the general protocol of the method, the specific protocol of the system and the run cards of the individual executions. The complete record is distributed together with the database as an RO-Crate. In the case of Cr--Ni, this record was sufficient to regenerate the database on a second machine, with the predictions of the model reproduced within the declared tolerances. Only the parameters along nearly flat directions of the objective differed between the reruns, and these parameters are reported in the record as undetermined by the evidence.

Beyond reproducibility, the standard is a first step towards CALPHAD assessments carried out by AI agents, since the expert knowledge that guides an assessment is encoded in protocols that a machine can follow. Since each protocol is written in a schema that a machine can read and each step is linked to the exact command that executes it, an agent with no prior knowledge of an assessment can read the protocol, run its steps and compare the results with the original ones. The rerun of the Cr--Ni assessment through the ORCHID-SAGE server provides direct evidence of this, since an agent with access only to the tools of the server and to read-only file tools regenerated the database without any other information and reproduced its predictions within the declared tolerances. The same server can also be used to start new assessments, in which the agent gathers the evidence from ORCHID and works through the general protocol step by step, although the fit steps for a new system are still adapted by a person, and the decisions that the protocol reserves for people are requested from them and recorded together with their reasons. ORCHID-SAGE thus provides a framework for AI-driven CALPHAD assessments that remain open to inspection, because the actions of the agents and the decisions of the people involved all become part of the record from which the database can be regenerated. Human guidance, however, remains essential in such assessments. Choosing phases and models, accepting new evidence, and adopting new rules require scientific judgment that agents do not replace, and the standard's role is to ensure this judgment is exercised explicitly and recorded, rather than left implicit in the final parameters.

The assessment itself, which combines the evidence of ORCHID with the one-shot model of PhaseForge and the gradient-based refinement of PhaseForge+, reproduces the measured eutectic and the top of CrNi$_2$ to within a few kelvin. It also points to the Cr-rich boundary of the bcc solid solution as the region of the diagram in which the model and the evidence still disagree, and in which new measurements would therefore be most useful. These results show that the reproducibility of a CALPHAD assessment is within reach at a modest cost, provided that the decisions of the assessor are recorded with the same care as the data. A database maintained in this way is no longer fixed at the time of its publication, and it can instead become a living database, refined as self-driving laboratories and other sources provide new evidence.
